\section*{Introduction}

Cerebrovascular disease represents one of the leading causes of death and long-term disability worldwide. In 2021, stroke alone accounted for an estimated 93.8 million prevalent cases and 7.3 million deaths, constituting over 10\% of global mortality \cite{wso2025,gbd2021stroke}. Beyond acute cerebrovascular events, a spectrum of chronic intracranial pathologies --- including unruptured intracranial aneurysms, arteriovenous malformations, and stenoocclusive disease --- carries substantial risk of catastrophic neurological outcomes. A critical, and still incompletely resolved, challenge in the management of these conditions is the reliable, non-invasive quantification of intracranial hemodynamics. Blood flow velocity, wall shear stress (WSS), and pressure distributions within the cerebral vasculature are increasingly recognised as key biomechanical determinants of lesion formation, growth, and rupture risk, yet they remain inaccessible to conventional clinical imaging \cite{wahlin2022,cebral2019wss}.

Four-dimensional Flow MRI (4D Flow MRI) has emerged as the principal non-invasive modality for comprehensive, time-resolved haemodynamic assessment of the cerebral circulation \cite{markl2012,dyverfeldt2015}. By encoding three-directional velocity information across the entire cerebrovascular tree within a single acquisition, 4D Flow MRI enables simultaneous evaluation of arterial and venous flow, net flow quantification, retrospective vessel-of-interest analysis, and the derivation of advanced haemodynamic metrics including WSS and oscillatory shear index \cite{bock2010,stienen2020wss}. Consequently, 4D Flow MRI is increasingly applied to a broad range of cerebrovascular conditions, from the haemodynamic characterisation of intracranial aneurysms and carotid bifurcation disease to the investigation of cerebral autoregulation impairment and flow alterations in neurodegenerative conditions \cite{wahlin2022,sun2020cfd}.

Despite its clinical promise, the practical utility of 4D Flow MRI at standard clinical field strengths (1.5T and 3T) is significantly constrained by its inherently limited signal-to-noise ratio (SNR), spatial resolution, and velocity-to-noise ratio (VNR). These limitations are particularly pronounced in the intracranial circulation, where vessel diameters are small, flow velocities are comparatively low, and accurate haemodynamic quantification demands sub-millimetre spatial resolution \cite{markl2012,dyverfeldt2015}. The partial-volume effect introduced by low spatial resolution systematically biases velocity and WSS estimates, underestimating peak flow velocities and distorting the spatial topology of the velocity field \cite{sun2020cfd}. Furthermore, the extended acquisition times required for whole-brain 4D Flow coverage at 3T compound physiological noise contributions and limit clinical feasibility in acutely ill patients.

Ultra-high-field MRI at 7 Tesla (7T) addresses many of these fundamental limitations. The approximately linear increase in SNR with field strength, combined with improved susceptibility contrast, yields substantially enhanced spatial resolution and vascular conspicuity compared to 3T \cite{kara2022_7t,smra2024}. In the context of 4D Flow MRI, these gains translate directly into higher VNR, reduced partial-volume contamination, and the ability to resolve smaller intracranial vessels --- enabling, for the first time, haemodynamic characterisation at the level of perforator arteries and the distal cortical circulation. Recent work has demonstrated that 7T 4D Flow MRI can achieve repeatable, quantitatively accurate intracranial flow measurements within clinically acceptable acquisition times \cite{mccormick2020_7t4dflow}. However, the limited global availability of 7T systems, the associated infrastructure and safety requirements, and the challenges of integrating investigational ultra-high-field acquisitions into routine clinical workflows place severe constraints on the direct clinical translation of 7T 4D Flow MRI.

This gap between the quality achievable at 7T and the practicality of routine 3T acquisition motivates a computational approach: training deep neural networks to learn the mapping from 3T 4D Flow MRI data to their 7T counterparts, thereby enabling virtual 7T quality from widely available clinical hardware. Deep learning has demonstrated remarkable success in MRI super-resolution and enhancement tasks, learning to recover fine structural detail from low-resolution or low-SNR inputs \cite{duffy2022sr,wang2024srrt}. However, existing super-resolution approaches for 4D Flow MRI are typically trained and validated on synthetic downsampling pipelines that do not capture the complex inter-scanner differences arising from distinct hardware, reconstruction algorithms, and k-space sampling patterns between field strengths \cite{ferdian2020,elahmar2023}. We address this by formulating two complementary learning tasks: a same-resolution domain translation task ($\times$1) that maps registered 3T acquisitions to their masked 7T counterparts, and a combined domain-translation and super-resolution task ($\times$2) in which spatially downsampled ($\times$0.5 trilinear zoom) 3T volumes are used as inputs and the full-resolution masked 7T volumes as targets. In both cases the ground truth is defined exclusively within the Circle of Willis mask $\mathcal{M}$, focusing supervision on the haemodynamically relevant intraluminal region. Bridging the real-world domain gap between 3T and 7T 4D Flow MRI requires paired, in-vivo acquisitions --- a prerequisite that has not previously been systematically addressed in the context of cerebrovascular 4D Flow.

Domain adaptation and cross-scanner generalisation are active areas of research in medical image analysis. Differences in MRI scanner platforms introduce systematic biases in image contrast, noise characteristics, and artefact patterns that degrade the performance of models trained on one domain when applied to another \cite{glocker2023domain,perone2019transductive}. In the specific context of multi-field-strength image synthesis, deep learning approaches have demonstrated the feasibility of learning 7T-like tissue contrast and spatial detail from 3T structural MRI \cite{qu2020synth7t,cui2024synth7t}, but no comparable work has addressed the fundamentally more complex case of quantitative, vector-valued 4D Flow MRI, where the network must predict physically consistent, three-directional velocity fields rather than scalar image intensities.

Physics-informed learning offers a principled framework to incorporate domain knowledge --- the governing equations of fluid mechanics and the physical constraints of MRI signal formation --- into the training objective, improving both the physical plausibility and generalisation of the learned haemodynamic representations \cite{kissas2020,ayed2023piml}. In the context of 4D Flow MRI, leveraging the divergence-free constraint of incompressible flow and the vessel-mask-aware structure of vascular velocity fields allows the network to focus its representational capacity on the haemodynamically relevant intraluminal regions, suppressing noise and artefacts in the extravascular background.

In this work, we present \textbf{NeuroFlow}, a supervised deep learning framework for joint 3T-to-7T domain translation and super-resolution of cerebrovascular 4D Flow MRI. The key contributions of this work are as follows:

\begin{enumerate}
    \item We acquire and process the first dataset of paired 3T/7T 4D Flow MRI studies in patients with cerebrovascular pathology, establishing a rigorous multi-stage registration and motion correction pipeline to achieve voxel-level correspondence between field strengths.

    \item We introduce a dual-branch vascular segmentation strategy combining a topology-aware nnU-Net ensemble \cite{isensee2021nnunet,yang2023topcow} with classical Frangi/Sato vesselness filtering \cite{frangi1998,sato1998} to generate high-quality Circle of Willis masks that guide mask-aware training.

    \item We adapt and extend the SR4DFlowNet architecture \cite{ferdian2020,elahmar2023} with attention mechanisms --- including CBAM \cite{woo2018cbam} and bidirectional cross-attention \cite{vaswani2017attention} --- and evaluate two complementary task configurations: (i) same-resolution 3T$\to$7T domain translation ($\times$1), mapping registered 3T inputs to their masked 7T counterparts; and (ii) joint domain translation and super-resolution ($\times$2), mapping trilinearly downsampled ($\times$0.5) 3T inputs to the full-resolution masked 7T target. In both tasks the ground truth is the registered 7T volume restricted to the CoW mask $\mathcal{M}$.

    \item We design a composite, mask-aware loss function that separately penalises velocity reconstruction error within and outside the vascular lumen, with an additional isotropic total variation penalty \cite{rudin1992tv} on extraluminal regions to suppress physiologically implausible background predictions.

    \item We demonstrate that the resulting framework achieves substantial reductions in relative speed error (RSE) and relative magnitude error (RME) within the Circle of Willis, representing a meaningful step towards virtual 7T quality cerebrovascular haemodynamics from routine 3T acquisitions.
\end{enumerate}

The remainder of this paper is organised as follows. Section~2 reviews related work on 4D Flow MRI processing, MRI super-resolution, and domain adaptation. Section~3 describes the methodology in detail, covering data acquisition, registration, network architecture, loss formulation, and training protocol. Section~4 presents quantitative and qualitative results across architectural variants and task configurations. Section~5 discusses the clinical implications, limitations, and directions for future work.

\bibliographystyle{ieeetr}
\bibliography{neuroflow}
